\section{Splitting the data}
\label{part:split}

Splitting decides what the final score actually measures, so both common ways of doing it were built and compared.

The \textbf{random split} divides the 1965 unique molecules into 70\,\% training, 15\,\% validation and 15\,\% test, keeping the proportion of positive labels the same in each part, with seed 42.

The \textbf{scaffold split} first reduces every molecule to its Bemis--Murcko scaffold \parencite{Bemis} --- the ring systems plus the pieces linking them, with all side chains stripped off. Molecules that share a scaffold are close relatives, often differing by nothing more than one added methyl group, so they are nearly the same molecule as far as a model is concerned. The split groups molecules by scaffold and fills training, then validation, then test, starting from the largest group and never breaking a group up. Two situations need a rule of their own:

\begin{itemize}
    \item \textbf{Molecules with no rings} have an empty scaffold. They are all treated as one group so that they end up together, because they have no ring system to generalise from and spreading them out would put the easiest cases in the test set.
    \item \textbf{Groups of equal size} are shuffled with a fixed seed before sorting. This is not cosmetic. Measured in blocks of 250 rows, the share of positive labels in the file runs 0.656, 0.524, 0.532, 0.428 and then exactly 1.000 for every block from row 1000 onwards --- the second half of the file is sorted by label. Breaking ties by position in the file therefore takes all the one-molecule scaffolds from that all-positive tail and produces a test set containing only one class, on which ROC-AUC cannot even be computed. Shuffling removes the problem and changes nothing else, since groups are still never divided.
\end{itemize}

\begin{table}[htb]
    \centering
    \caption{The two splits over the 1965 unique molecules; the overall share of positives is 0.7634. Neither split shares a molecule between subsets; the random split shares \textbf{238 scaffolds}, the scaffold split none.}
    \label{tab:splits}
    \footnotesize
    \begin{tabular}{@{}llrrr@{}}
        \toprule
        Split & Subset & $n$ & Pos.\ rate & Scaffolds \\
        \midrule
        Random   & train & 1375 & 0.7636 & 773 \\
                 & val   &  295 & 0.7627 & 225 \\
                 & test  &  295 & 0.7627 & 226 \\
        \midrule
        Scaffold & train & 1375 & 0.7665 & 430 \\
                 & val   &  294 & 0.7585 & 294 \\
                 & test  &  296 & 0.7534 & 296 \\
        \bottomrule
    \end{tabular}
\end{table}

\Cref{tab:splits} shows both. The sizes and the share of positive labels agree with the full dataset to within half a percentage point, and both splits cover all 1965 molecules without overlap. They differ in exactly one number: in the random split, 238 scaffolds appear on both sides of the divide, while in the scaffold split none do.

That single number is the point of the exercise. When a test molecule shares its scaffold with a training molecule, the model has already seen a close relative with a measured label, and it can score well by recognising the family rather than by understanding anything. Real use is the opposite situation --- a chemist asks about a new kind of molecule nobody has measured --- and the scaffold split is built to imitate that. The difference between the two scores is therefore not a flaw to be fixed; it is an estimate of how much of the random-split score was recognition.

Both test sets were then set aside and not looked at again until the single final evaluation in \ref{part:eval}.
